\section*{Aktivitas Pembelajaran dan Praktikum}
\addcontentsline{toc}{section}{Aktivitas Pembelajaran dan Praktikum}

\begin{table}[H]
\centering
\caption{Rencana aktivitas Bab VII}
\label{tab:ch07-aktivitas}
\begin{tabularx}{\textwidth}{@{}llX@{}}
\toprule
\textbf{Sesi} & \textbf{Bentuk} & \textbf{Kegiatan} \\
\midrule
1 & Ceramah & Pola scanner: \code{start}, \code{current}, dan lookahead \\
2 & Live coding & Lexer tangan: identifier, angka, operator \\
3 & Live coding & Lexer Flex dan perbandingan aturan \\
4 & Praktikum & M1: implementasi, pengujian, dan laporan \\
\bottomrule
\end{tabularx}
\end{table}

\subsection*{M1 / Praktikum 1: Lexer}
\textbf{Tujuan.} Menghasilkan lexer Subset-C yang bekerja, diuji, dan
terdokumentasi.

\textbf{Keluaran.}
\begin{enumerate}[label=\arabic*.,leftmargin=*]
  \item Lexer tangan dalam C (\file{lexer.c}) yang lulus tes.
  \item Lexer Flex (\file{lexer\_flex.l}) yang menghasilkan token
        setara.
  \item Minimal 30 kasus uji dengan berkas harapan.
  \item Keluaran program: aliran token kedua lexer pada berkas
        masukan uji.
  \item Laporan singkat: desain, hasil tes, kendala, dan refleksi.
\end{enumerate}

\textbf{Langkah kerja.}
\begin{enumerate}[label=\arabic*.,leftmargin=*]
  \item Lengkapi \file{include/token.h} dan \file{include/lexer.h}.
  \item Implementasikan \file{src/lexer.c} dan \file{src/error.c}.
  \item Implementasikan \file{src/main\_lex.c}.
  \item Implementasikan \file{src/lexer\_flex.l}.
  \item Jalankan \file{build.bat} dan \file{run-tests.bat}.
  \item Bandingkan keluaran lexer tangan dan Flex.
  \item Tulis laporan dan serahkan.
\end{enumerate}

\textbf{Rubrik.}
\begin{table}[H]
\centering
\caption{Rubrik Praktikum 1}
\label{tab:ch07-rubrik}
\begin{tabularx}{\textwidth}{@{}lXr@{}}
\toprule
\textbf{Aspek} & \textbf{Indikator} & \textbf{Bobot} \\
\midrule
Kebenaran & Semua token dikenali, tes lulus & 40\% \\
Kualitas kode & Struktur, penamaan, komentar, tanpa peringatan & 30\% \\
Kelengkapan & Tangan dan Flex, 30 tes, laporan & 20\% \\
Pengujian & Kasus normal, batas, dan kesalahan & 10\% \\
\bottomrule
\end{tabularx}
\end{table}

\textbf{Definisi selesai.}
\begin{checklist}
  \item Lexer tangan dan Flex berjalan.
  \item Minimal 30 tes lulus.
  \item Keluaran kedua lexer identik pada berkas uji.
  \item Diagnostik menyebut lokasi dan penyebab.
  \item Laporan diserahkan.
\end{checklist}
